\scanpage{5}
\uline{Proposition 1} Les faisceaux $\underline L^{\otimes n}$ sont coh. plats en ttes dimensions. Le faisceau $f_*(\underline L^{\otimes n})$ est localement libre, de rang $0$ si $\unsure{n|p}<0$, $1$ si $n=0$ ou $1$, $n\deleted{\unsure{-1}}$ si $n\geq1$. \deleted{\unsure{Le|$\underline L$}} Le faisceau $\underline L$ est ample, et de façon précise $\underline L^{\otimes n}$ est \deleted{\uline{très}} \uline{ample} \deleted{pour}\inserted{ssi} $n\geq1$, très ample ssi $n\geq3$

La formule de RR, où $g=1$, d'où $2g-2=0$, nous montre que \uline{tout faisceau} \deleted{loc libre} inversible sur $C$, de degré \inserted{relatif} $d\ne0$, est coh. plat, et que sa dimension est \deleted{\unsure{donné|donc}} nulle si $d<0$, et $d\deleted{\unsure{-(g-1)}}$ si $d\geq1$. Donc la première assertion est démontrée pour $n\ne0$, et pour $n=0$ elle est connue par ailleurs, car $\underline L^{\otimes n}=\underline O_C$.
% slanted
\margin{prop. préliminaire sur un faisceau inversible quelconque}

\inserted{On \deleted{\unsure{sait}} sait qu'un faisceau inv. sur $C$ de degré relatif $d$ est ample ssi $d\geq0$, donc $\underline L^{\otimes n}$ est ample ssi $n\geq1$.} \deleted{Et puisque le \unsure{$\underline L$ est de degré relatif $1>0$|$\underline L$ est de degré relatif $d>0$}, il est ample, pour voir s'il \unsure{donc} $\underline L^{\otimes n}$ est très ample pour $n>0$.}

Si $\underline L^{\otimes n}$ est très ample, on a $n\geq1$, d'ailleurs on ne peut avoir $n=1,2$ car alors $\mathbb P(f_*(L^{\otimes n}))$ est de dim relative $0$ ou $1$, donc $C$ ne peut s'y immerger (car il y serait isomorphe, absurde), donc il faut $n\geq3$. Celà suffit, car à cause de la platitude cohomologique, \deleted{et} \deleted{\unsure{est}} on est ramené au cas \deleted{où} $S=\unsure{\mathrm{Spec}|{\mathrm{Sp}}}\,k$, $k$ un corps\inserted{ alg. clos}. \deleted{Mais alors D'ailleurs, comme $\underline L$ est engendré loc$^{\rm t}$ par ses sections} Mais \deleted{alors} comme $3>(2g-2)+2=2g=2$, un raisonnement connu n'utilisant que RR donne le résultat.
% archive sheet number: 4
